\begin{frame}[fragile]{Symmetric Encryption}

\begin{itemize}
	\item Encryption is the process of making a message confidential (like a hash) while allowing it to be reversible (decrypted) with the proper key. 
	\item In symmetric encryption, the same key is used to encrypt and decrypt the message.
\end{itemize}

\begin{center}
\includegraphics[width=0.5\textwidth]{topics/SE_tools2/security/pic/symmetry/sym.png}
\end{center} 


\end{frame}


\begin{frame}[fragile]{Node.js module}

\begin{itemize}
	\item We use the `crypto' module to use JavaScript decrypt and encrypt functions. 
	\item Each time a message is encrypted, it is randomized to produce a different output. So, we should also use the random function. 
\end{itemize}

\begin{Code}{}%
\begin{lstlisting}[firstnumber=1, label=glabels, xleftmargin=0pt, basicstyle=\scriptsize]% 

const { createCipheriv, randomBytes, createDecipheriv } = require('crypto');

/// Cipher
const message = 'I like ASE 285 students!';
\end{lstlisting}%   
\end{Code}%

\end{frame}

\begin{frame}[fragile]{}

\begin{itemize}
	\item The first step is to create a key and an initialization vector (iv) as in lines 2 -- 3. This information becomes the key that is shared. 
	\item The next step is to create a cipher object as in line 5.
\end{itemize}

\begin{Code}{}%
\begin{lstlisting}[firstnumber=1, label=glabels, xleftmargin=0pt, basicstyle=\scriptsize]% 

// these key/iv should be shared. 
const key = randomBytes(32); const iv = randomBytes(16);
const infoToShare = {key:key.toString('hex'), iv:iv.toString('hex')}

const cipher = createCipheriv('aes256', key, iv);
\end{lstlisting}% 
\end{Code}%

\end{frame}

\begin{frame}[fragile]{Encryption Code}

\begin{itemize}
	\item We can use the cipher object to encrypt the message. 
	\item We should add the final code to make it harder to decrypt. 
\end{itemize}

\begin{Code}{}%
\begin{lstlisting}[firstnumber=1, label=glabels, xleftmargin=0pt, basicstyle=\scriptsize]% 

const encryptedMessage = cipher.update(message, 'utf8', 'hex') + cipher.final('hex'); 
\end{lstlisting}%   
\end{Code}%

\end{frame}

\begin{frame}[fragile]{Decryption Code}
\begin{itemize}
	\item Once we have the shared key and iv, we can decrypt the message. 
\end{itemize}

\begin{Code}{}%
\begin{lstlisting}[firstnumber=1, label=glabels, xleftmargin=0pt, basicstyle=\scriptsize]% 

const key2 = Buffer.from(infoToShare.key, 'hex')
const iv2 = Buffer.from(infoToShare.iv, 'hex')
const decipher = createDecipheriv('aes256', key2, iv2);
const decryptedMessage = decipher.update(encryptedMessage, 'hex', 'utf-8') + decipher.final('utf8');
\end{lstlisting}%   
\end{Code}%
\end{frame}

\begin{frame}[fragile]{AES 256 Algorithm}

\begin{itemize}
	\item In the examples, we have used AES (Advanced Encryption Standard) 256 algorithm to encrypt information. 
	\item This is a standard published and maintained by the NIST. 
	\item We should never come up with our own hidden encryption/decryption algorithms, but we should use the standard algorithms such as AES 256. 
\end{itemize}
\end{frame}

\begin{frame}[fragile]{}

\begin{itemize}
	\item It is a well-known fact that the AES algorithm cannot yet be cracked, at least not in this lifetime. 
	\item It would take billions of years for a supercomputer to crack even a 128-bit AES key. 
	\item Quantum computers can break AES algorithms quicker, but, according to some sources, it would still take a quantum computer roughly six months to exhaust the possibilities of a 128-bit AES key.
	\item So, we can say that it is safe to use AES 256 in our applications. 
\end{itemize}
\end{frame}

\begin{frame}[fragile]{Can we ...}

\begin{itemize}
	\item Can we share keys without exchanging the keys?
	\item As exchanging keys might impact security, people have thought about sharing keys without exchanging them. 
	\item DH Key Exchange Algorithm is the most widely used method to share keys without actually exchanging them. 
\end{itemize}

\end{frame}

\begin{frame}[fragile]{DH (Diffie-Hellman) Key Exchange}

\begin{columns}%
\begin{column}{0.7\textwidth}%
\begin{itemize}
	\item  DH (Diffie-Hellman) is a method of securely exchanging cryptographic keys over a public channel.
	\item Alice and Bob has random integer a and b that they do not share, but they can create $A = g^a mod\ p$ and share p, g, and A with Bob. 
	\item Then Bob creates $B = g^b mode\ p$ to send B to Alice. 
\end{itemize}
\end{column}%
\begin{column}{0.3\textwidth}%
\begin{center}
\includegraphics[width=0.99\textwidth]{topics/SE_tools2/security/pic/dh.png}
\end{center}
\end{column}%
\end{columns}%

\end{frame}

\begin{frame}[fragile]{}

\begin{itemize}
	\item Hackers can steal $p,g,A,B$, but not a (private only to Alice) and b (private only to Bob).
	\item Both Bob and Alice can generate a key $g^{ab} mod\ p$ from $(g^a)^b mod\ p$ = $(g^b)^a mod\ p$ = $g^{ab} mod\ p$, but not the hacker because the hacker cannot make $g^{ab} mod\ p$ from $p, g, A, B$. 
	\item So both Bob and Alice can share the key with this simple integer arithmetic. 
\end{itemize}

\end{frame}